\chapter[Why May the State Command Us?]{What Gives the State the Right to Command People?}
\section{Why Can a Coin Buy a Steamed Bun?}
First, pick up an object: a one-yuan coin.

Weigh it in your palm. It is light, made of stainless steel. As a piece of metal, it may be worth only a few fen. Yet take it into a breakfast shop and you can exchange it for a steaming bun. The shopkeeper accepts it without blinking, as though receiving something self-evidently valid.

Now inspect the design: one face shows "1 yuan" and a chrysanthemum, the other the national emblem of the People's Republic of China. What makes this metal money is not the metal but the promise behind that emblem. Something called the state guarantees that it works and decrees that nobody may refuse or counterfeit it. Counterfeiting a banknote uses only paper and ink at negligible cost, but criminal law and prison await you.

Think also of the traffic lights on your way to school. A red light appears, and dozens of cars, each weighing a tonne or two and capable of easily crossing the white line, stop together. There is no wall or ditch at the intersection, often not even a traffic officer. What stops the steel is one light and the whole arrangement behind it: laws, driving licenses, fines, cameras, and the thought in everyone's mind that a red light means one ought to stop.

Here is our question. An enormous entity can order you about: the age you must attend school, the minimum age for cycling, how much tax you owe on earnings, what you may not say or touch. Disobey and it has police, courts, and prisons. But what gives it that right? A robber forces you at knifepoint to surrender your wallet. You comply, but nobody thinks the robber is entitled to your money. What distinguishes the state from the robber: greater violence, or something deeper?

We will find that even this enormous entity's name requires careful distinctions. First, come to a scene where it is absent and see what life looks like when it disappears.

\section{Life Without Leviathan}
Time: autumn, 1644. Place: an ordinary village in central England.

King Charles I and Parliament have been fighting a civil war for two years. The king says God appointed him monarch and everyone must obey. Parliament says even a king cannot arbitrarily tax or arrest people: power needs boundaries. Neither persuades the other, so both take up swords.

Imagine being a miller's son in the village. Early in the morning you smell smoke, not cooking smoke but smoke from the neighboring village to the east. A cavalry unit passed yesterday, requisitioning grain "for His Majesty." They took the food and burned two barns whose owners refused. Nobody can tell whether they were royal or parliamentary troops, and nobody cares: both sides take grain and claim to represent justice.

Your father hides the family's only two pigs in the cellar. Several checkpoints stand on the road, each demanding a toll. The collectors change flags and collect again. School stopped long ago. People dare not hold markets because gathering may lead to arrest for unlawful assembly. Some of the village's strongest young men were taken by royal troops; others joined Parliament's army; others disappeared into the woods as bandits. When everyone carries guns, the boundary between soldier and robber is already blurred.

At night you cannot sleep, not from fear of a particular person but because you cannot know who will knock tomorrow, under which flag, applying which rules. The worst torment is not bad rules but no rules at all.

Under the same sky, an Englishman named Thomas Hobbes lives in exile in Paris. He has escaped the fighting, but not the question in his mind. In 1651 he publishes a book named after an enormous sea monster: Leviathan. It depicts life without a state. Without a common power restraining everyone, people confront one another like wolves. Nobody can farm, work, or sleep securely, since someone stronger may always rob them. His conclusion unsettled many readers then and now: to escape that horror, people would rather create a powerful Leviathan and give it authority to command everyone.

Remember this village. Most of our debate begins there. One side says life without a state looks like this, so the state has the right to command. The other replies: must escaping fear of bandits mean fearing the king forever?

\section{By What Right? A Genuine Question}
First pin down the question itself, before hearing answers.

Every day you obey instructions. The student on classroom duty says it is your turn to clean the board, and you clean it. A red light appears, and you stop. A teacher forbids talking in class, and you fall silent. Now a seemingly argumentative question: what exactly are you obeying?

Notice three very different things. First, \textbf{force}: a robber's knife or an older child's fists. This obedience contains no ought, only must. Second, \textbf{reasons}: a classmate explains a problem and you follow their reasoning. This is not really obedience to commands but persuasion. The third is distinctive: \textbf{authority}. The duty student tells you to clean the board not because they can beat you, perhaps they cannot, nor because they have proved it needs cleaning today, which they have not. They tell you because they hold that role. Their position turns their words into a command.

Political philosophy gives this third kind of obedience a name: \textbf{legitimacy}, also understood as justified authority. Take an everyday example. "Clean the board" is an exercise of duty from the duty student but interference from a passing student in the next class. The difference is not the instruction's content but whether the speaker is \textbf{entitled} to issue it, and where that entitlement comes from. Power possessing force alone is an oversized robber. With legitimacy, those obeying feel not that they are being robbed but that they are fulfilling an obligation.

The real question now appears as four linked questions:

\begin{itemize}
\item What distinguishes a state's command from a robber's? This is the question of legitimacy.
\item If there is a difference, who grants it? Have those commanded \textbf{consented}? This is the question of consent.
\item Does majority approval permit doing anything to a minority? This is the question of boundaries.
\item If the law itself is wrong, should one obey or resist? This is the question of conscience.
\end{itemize}
Humanity has seriously argued each question for centuries. This chapter will not settle them for you, only lay every side's cards on the table. Before continuing, remember your intuition now. Is stopping at a red light really different from being stopped by a robber? If so, how?

\section{What Order's Defenders and Skeptics Fear}
Now hear the different voices around this question. We will first state each case fully, without judging it.

\textbf{Voice one: defenders of order fear disorder more than tyranny.}

Give their case its strongest form. They are often dismissed as simply wanting someone to control them, which is unfair. Remember the village in 1644. Their argument has at least three layers.

First, security is the prerequisite for everything. You may complain about heavy taxes and excessive rules, but complaining requires being alive, having an unburned shop, and daring to leave home tomorrow. Freedom, dignity, and rights all sit inside the box of order. Break the box and nothing inside remains safe. Order is therefore not one benefit among many but the container for all benefits.

Second, ordinary people calculate this too. Support for a strong state is not merely rulers' propaganda. Throughout history, farmers, stallholders, and millers have often sincerely chosen harsh government over chaos. Not because they lack self-respect, but because they have paid disorder's bill: confiscated grain, missing sons, unopened markets. When a grandmother who survived war tells her grandson complaining about queues that having someone in charge is already a blessing, a whole history of blood and tears lies behind her words.

Third, concentrated power offers real efficiency. Building embankments, preventing epidemics, digging wells, making roads, and maintaining currency require decisions, coordination, and compulsory cost-sharing. A society where nobody may command anyone might fail to build even a decent drain.

This is a substantial case. It underpins many ordinary comforts you enjoy: going out at night, trusting mobile payments, and receiving vaccines.

\textbf{Voice two: skeptics fear uncaged power more than disorder.}

But skeptics also have a complete argument, drawn from equally bloody experience.

First, Leviathan can rob too. In the village of 1644, bandits were never the only grain thieves. Royal and parliamentary troops were equally accomplished. Giving one organization a monopoly of violence requires assurance that it behaves better than scattered bandits. Who guarantees that? Bandits leave after robbing you; a robbing state wears uniforms, carries documents, and returns every year. Is one certain large robber really a good bargain against several possible small ones?

Second, consent is often fictitious. The state says you consented to its rule. When? You signed no contract, merely happened to be born here. Does birth count as a signature? For most people, skeptics argue, consent to government resembles consent to the weather: it is a situation, not a choice. Calling that situation a voluntary contract makes the commanded cover for their commanders' falsehood.

Third, order's benefits are distributed very unevenly. The same order gives some protection and others restraint. Among people afraid to go out at night, some fear criminals, others precisely those patrolling the streets. Whether the state has made you safe depends on who you are.

\textbf{A third voice comes from China two thousand years ago.}

This debate is no Western monopoly. During the Warring States period, Mencius made a statement that unsettled successive emperors. Mencius, "Fathoming the Mind, Part Two," records:

\begin{quote}
The people matter most, the altars of soil and grain come next, and the ruler matters least.
\end{quote}
Soil and grain refer to their deities and stand for state power. The statement ranks the people first, the state next, and the monarch last. Notice how it reverses the defenders' claim that monarch and state are the container: the container serves what is inside, not the reverse. Later emperors repeatedly sought to remove or alter this sentence, showing how precisely it struck the nerve. Accepting that order means accepting conditional, revocable power.

Incidentally, ancient Chinese Daoism offers a more radical voice: the best government leaves people barely aware of a ruler. The Daodejing says, broadly, that after the finest ruler completes the work, people say they have always lived this way themselves. This may be among the earliest Eastern echoes of government governing least.

Set the sides together and an honest fact emerges. They fear different things, and both fears have bodies as evidence. Disorder and tyranny have each killed millions. The difficulty is not choosing a side but acknowledging both fears as real.

\section[State, Nation, and Government]{Thinking Bridge: Separate State, Nation, and Government}
Before comparing theories, sharpen a small tool. It can cut through at least half the confused arguments in news, online, and over dinner.

The tool separates three boxes: \textbf{state, nation, and government are not the same thing}.

\textbf{State}: a political machine with defined boundaries, territory, a permanent population, institutions monopolizing violence, including armed forces, police, courts, and taxation, and recognition by other such machines. Its keyword is \textbf{institutions}. The United Nations assigns seats by state, regardless of who has just taken office.

\textbf{Government}: the people and administration currently operating that machine. Its keyword is \textbf{tenure}: governments change, are replaced, or are dismissed. The state is the vehicle; government is the driver. A new driver leaves the same vehicle. Criticizing bad driving does not mean wanting to smash the car, still less failing to care about it.

\textbf{Nation}: people who regard themselves as an us, a community joined by shared language, history, memory, culture, and sometimes religion. Its keyword is \textbf{identity}. A nation is an imagined community, not because it is false but because it lives in mutual identification rather than boundary stones and tax offices.

Separate these words and many knots immediately loosen:

\begin{itemize}
\item \textbf{A nation without a state}: Kurds live around the borders of four Middle Eastern countries, with shared language and historical memory but no state machine called Kurdistan. Their relationship to states differs completely from the textbook's default picture. Their us corresponds to no passport.
\item \textbf{One state containing several nations}: Switzerland uses four official languages. The ancestors of its German-, French-, and Italian-speaking populations were not originally one group, yet they have lived within one state machine for centuries. Nation and state are not naturally conjoined twins.
\item \textbf{A new government while the state remains}: Britain changes prime ministers like seasonal shop displays. Its state institutions maintain continuity between Elizabethan times and today, while governments have changed hundreds of times. Opposing this government and opposing this state are fundamentally different. You can now recognize deliberate online conflation of the two.
\item \textbf{Identities may overlap or conflict}: someone may identify simultaneously with nation and state, or feel the two clash. Many historical wars and divisions arose from mismatched national and state boundaries.
\end{itemize}
Apply the tool whenever a statement uses these words. Does it mean the machine, the driver, or us? Which appears in winning honor for one's country, changing government, or national feeling? Without these distinctions, anger easily strikes the wrong target, and others can easily steer it there.

\section{Three Imaginary Contracts}
Now read a short primary source, actually three. The skeptic asked a sharp question: when did you consent to state rule? European thinkers of the seventeenth and eighteenth centuries faced it directly. Their method was a thought experiment. \textbf{Suppose} people once lived without a state, then negotiated a contract to establish one. What terms would rational people sign? This is the famous social-contract approach. The contract is imaginary but serves a real purpose: it measures the state's right to command by whether its powers meet the hypothetical agreement's terms.

The problem is that the three best-known contract designers drew up entirely different terms.

\textbf{First: Hobbes's survival contract.} As noted, Hobbes had just escaped civil-war England. Chapter 13 of Leviathan describes the absence of common power: broadly, no agriculture, commerce, building, or knowledge, because with life liable to violent theft, who would invest in anything long-term? Worst of all is perpetual fear of violent death. The passage ends with one of political philosophy's most quoted descriptions: life in that condition is solitary, poor, degraded, brutal, and brief.

Hobbes's terms therefore require everyone to surrender all rights to use violence to one sovereign, either a monarch or an assembly, whose commands become law. In return come security and order. The key clause: \textbf{this contract is almost impossible to cancel}. Allowing easy withdrawal would return everyone to mutual hostility. Freedom here chiefly means what the sovereign has not explicitly forbidden, not a right to resist the sovereign.

\textbf{Second: Locke's property-protection contract.} John Locke wrote Two Treatises of Government around the Glorious Revolution, when England had actually expelled a king with little bloodshed. His state of nature is less grim than Hobbes's. People are generally rational and possess property but lack an accepted judge for disputes. If a neighbor occupies your apple orchard, you must argue your own case with a stick, risking escalation. Locke's terms are much looser: people surrender only their right to judge and punish others themselves, receiving governmental protection of life, liberty, and property. The key clause: \textbf{government is a trustee and can be replaced if it fails}. A ruler violating people's property and liberty breaches the agreement first, entitling them to replace him. North American colonists later copied this page directly into their declaration against Britain's king.

\textbf{Third: Rousseau's general-will contract.} Jean-Jacques Rousseau of Geneva lived later and faced a different problem: not chaos, but a society where wealth accumulated among the rich and every rule favored them. The opening sentence of The Social Contract still appears on countless French notebooks and posters:

\begin{quote}
Human beings are born free, yet everywhere they are in chains.
\end{quote}
Rousseau asks this chapter's central question: how can the chains be legitimate? His answer is the most radical. People do not hand authority to a sovereign; they \textbf{jointly} constitute the sovereign. Each obeys not another person's will but the general will formed by all, directed toward the community's common interest. Following it is not being commanded, because we made the law ourselves. Obeying it is obeying ourselves, which alone is genuine freedom.

Placed together, the contracts differ clearly. \textbf{What is surrendered}? Hobbes gives up everything, Locke only private adjudication, Rousseau gives it to us rather than anyone else. \textbf{What is received}? Survival, protected property, or obedience understood as freedom. \textbf{Can the agreement end}? Almost never; replacement after breach; sovereignty forever remaining with the people. One thought-experiment technique produces three states. The tool does not supply the answer; the premises inserted into the contract do. Its shape depends on how dangerous you assume human nature, how sacred property, and how reliable us to be.

\section{Three Explanations of Obedience}
Step back to a simpler fact: most people stop at red lights. Why? What happens is uncontroversial, but why it happens has at least three competing explanations. Notice what each explains and what escapes it.

\textbf{Explanation one: calculation, or obedience because it pays.} Running red lights risks fines, license points, and collisions. Unpaid taxes bring recovery proceedings and blacklisting. People weigh obedience's benefits against its costs and comply when benefits win. This neat account explains why weaker penalties or broken cameras quickly increase queue-jumping and illegal parking. But it cannot explain \textbf{obedience when nobody watches}. People still stop at empty intersections late at night and put rubbish in bins in camera-free alleys. If everything is calculation, what do they gain?

\textbf{Explanation two: habit and inertia, or obedience because things have always been this way.} Eighteenth-century Scottish philosopher David Hume made a cool observation: even powerful governments cannot subdue vast populations through force alone, since police and soldiers are always greatly outnumbered. Government rests on opinion: people are accustomed to obedience, believe they ought to obey, and see others obeying. Compliance resembles a riverbed worn by daily currents, not a contract signed one day. This fills calculation's gap: at the late-night red light, the body probably acts before the internal accountant. But it makes obedience too automatic to explain \textbf{why it suddenly collapses}. Whole generations have ceased obeying overnight and taken to the streets. Riverbeds change course; habit cannot explain the moment of change.

\textbf{Explanation three: belief, or obedience because authority seems entitled to it.} Early-twentieth-century German sociologist Max Weber offered a third account. People often obey because they sincerely regard authority as \textbf{legitimate}. He distinguished three sources: traditional authority, because ancestors have always done it this way, supporting patriarchs and hereditary monarchs; charismatic authority, because an exceptional leader's character and achievements inspire followers, supporting revolutionaries and founders; and legal-rational authority, because established procedures prescribe it. Here people obey rules rather than individuals. A teacher may grade you because the institution grants that power to the teaching role. A replacement teacher receives the same authority. Modern states chiefly rely on this third type. It accommodates what earlier explanations miss: late-night stopping because legal-rational belief has been internalized, and sudden collapse when evaporating belief can bring down a regime overnight. Regimes have often fallen with startling speed not because their armies vanished but because nobody believed in them anymore. Its limitation is excessive attention to belief at the expense of underlying interests and fears. Some profess belief because they have no choice, while universal belief itself may be the product of a propaganda machine working around the clock.

These explanations do not ask you to keep one and discard two. Your braking at a red light may combine calculation, habit, and belief. That reminder matters: a theory claiming to explain all obedience alone has probably quietly discarded something.

\section{Further Challenge: Is an Unjust Law Law?}
Now the chapter's weightiest question arrives. Everything so far assumes good laws that the state merely enforces. What if the law itself is wrong?

First encounter an ancient, famous counterexample correcting a common mistake: whatever the majority decides is legitimate.

Athens, 399 BCE. The city-state prides itself on democracy. Citizens vote on major matters, and courts have juries of hundreds of ordinary citizens selected by lot. That year, seventy-year-old philosopher Socrates faces charges of disbelieving in the city's gods and corrupting the young. The procedure is entirely lawful: accusation, public trial, arguments from both sides, then a vote by around five hundred jurors. Verdict: guilty. Another vote determines the sentence: death. Socrates refuses an escape arranged by friends. His student Plato later records the discussion in Crito. Broadly, Socrates argues that enjoying the city's laws throughout life amounted to tacit agreement. Escaping when the verdict goes against him would betray that agreement. He drinks poison and obeys the judgment.

Notice the entangled layers. \textbf{The procedure was lawful}: Athenians followed their rules. \textbf{The majority agreed}: both votes produced majorities. \textbf{Almost all later observers regard it as a miscarriage of justice}: Athens executed one of its greatest minds. Majority, procedure, and legality together still produced a wrongful judgment remembered for twenty-four hundred years. This is among the oldest evidence that majority decisions need boundaries protecting rights. Some things, such as a person's right not to be executed for ideas, should not enter the ballot box. Majorities can choose noodles or rice for lunch, not who disappears.

Socrates obeyed, but history offers another choice: \textbf{civil disobedience}, publicly and nonviolently violating a law one considers unjust, accepting punishment rather than evading it, and using that punishment to expose the law's problem to society.

Nineteenth-century American Henry David Thoreau was an early traveler on this road. Opposing slavery and the war against Mexico, he refused the poll tax and spent a night in jail. His later essay Civil Disobedience opens with a then-popular maxim: the best government governs least. He distinguished respecting law from outsourcing conscience to it. When law and conscience clash, one should first be upright, then law-abiding.

The road later extended farther. In 1930, Gandhi protested British colonial India's salt monopoly by leading followers nearly four hundred kilometers to the coast and openly making salt by boiling. Salt sustained poor people's lives. Was making it a crime? Thousands followed him in breaking the law and being arrested. Prisons filled, exposing colonial law's absurdity worldwide. More than twenty years later, in 1963, American minister Martin Luther King Jr. was arrested during protests against racial segregation. From Birmingham jail, he wrote a long letter answering white ministers urging lawfulness and patience. Citing the early Christian thinker Augustine, he argued, broadly, that an unjust law does not deserve the name law. He also explained civil disobedience's complete logic: willingness to accept punishment demonstrates respect for the rule of law itself. The target is this unjust law, not law's existence.

Behind these actions lies an unresolved legal debate. The \textbf{natural-law} tradition places a higher measure above written law: human dignity and basic justice. Laws failing it are merely commands dressed in judicial robes. \textbf{Legal positivism} holds that law consists of what is enacted through proper procedures; morality belongs to a separate account. Otherwise everyone makes their conscience a supreme court and the legal system disintegrates. Both have real arguments. Without natural law's measure, there would be no vocabulary to condemn Nazi Germany's lawfully enacted atrocities. The postwar Nuremberg trials established that merely following lawful orders cannot excuse crimes against humanity. Yet without positivism's stability, society risks the chaos of everyone establishing their own court.

This is what deeper inquiry reveals: freedom, authority, law, and conscience cannot all be maximized simultaneously. Mature civic thought does not choose one and discard three. It learns to recognize which two are pressing against one another where you stand.

\section[Every Agree Button You Have Clicked]{Back to Today: Every Agree Button You Have Clicked}
This centuries-old debate now inhabits daily life.

First, the classroom. A meeting votes on headphones during independent study: thirty-seven against, eleven for. A new prohibition enters the rules. This is miniature majority rule. It is broadly legitimate: everyone is affected, everyone has a vote, and those who lost can raise it again tomorrow. But suppose someone proposes voting to forbid Xiaoming ever speaking in class because he is noisy. Does voting still seem appropriate? Probably not. You sense that some things cannot be submitted to a vote: the basic right to speak, dignity against group humiliation, and the right not to have one's name blacked out. You have intuitively reached the rights-based limits on majority decisions. Democratic constitutions protect basic rights from fluctuating vote counts precisely to prevent society-wide versions of a class voting to isolate one person.

Now your phone. Every app registration produces an agreement tens of thousands of characters long and a button saying you have read and agreed. You have not read it; nobody has. But you click. Social-contract reasoning sees an agreement exchanging data for service. Is this the consent Hobbes, Locke, and Rousseau envisioned? They never faced a contract where refusal meant losing access to homework submission, group chats, and classmates. When your whole social life depends on clicking agree, how far is that consent from consent before a robber? No standard answer exists, but contemporary law has begun responding. Some countries rule that even lengthy agreements cannot remove certain rights, such as seeking accountability for data misuse. Some things above the agreement cannot be signed away. This is the newest version of the ancient claim that a higher measure stands above law.

Now online disputes. Someone criticizes a policy, and comments immediately accuse them of not loving their country. You now have a tool. The criticism targets \textbf{government}, the current operators and their decisions; the alleged victim is the \textbf{state}, the machine itself, or the \textbf{nation}, the community of us. Three concepts have been substituted. Translating criticism of driving into hatred of the entire vehicle and all passengers is a classic way to win arguments through confusion. Similarly, a majority condemning someone does not prove the condemnation just. Socrates's jury also had a majority. Numbers have never been justice's scale.

Finally, a closer issue. School requires phone inspections at the gate, which feel intrusive. The school invokes order: students cheat, photograph classmates, and play games in lessons. Defenders of order and skeptics confront each other at your gate, both reciting this chapter's lines. One says freedom needs order's container and one phone can destroy a classroom's concentration. The other asks where inspections without boundaries will stop: phones today, what tomorrow? This dispute persists because it is not misunderstanding but a real conflict between real values. You can do more than take sides. Ask about the rule's \textbf{source of legitimacy}: who made it, through what process, and with what channels for students and parents to speak? Ask about its \textbf{boundaries}: how far may inspection go? Ask about \textbf{remedies}: where can someone appeal mistreatment? These three questions translate the angry "What right do you have to control me?" into civic language.

\section{Your Question: Red Lights and Conscience}
Return to the opening intersection.

The light is red. No cars, officers, or cameras are present. You stop, perhaps from calculation, just in case; perhaps habit, as your foot stops itself; perhaps belief that rules are rules. One day the same light commands something different. It no longer merely asks you to stop, but to do something you sincerely regard as wrong.

Then you hold all this chapter's cards. Defenders of order ask you to remember the container and its contents: those who have paid disorder's bill know its cost. Skeptics say conscience cannot be outsourced wholesale, with Thoreau's and Gandhi's prison cells as evidence. Socrates stands between them as both warning, since majorities err, and puzzle, since he obeyed to the end. Social contract offers a ruler: does this command resemble terms rational people would accept? Conceptual distinction offers a mirror: do you oppose this rule, this government, or the whole us mistakenly made a target?

Answer, then: \textbf{under what conditions does a law lose its entitlement to your obedience}? When its procedure is unlawful? When the majority disagrees? When it violates rights that should never be voted on? Or when it crosses your inner boundary? If everyone draws that boundary differently, what should society do?

Give your answer and reasons. There is no standard answer, but one requirement: leave at least one seat in your reasoning for the other side's fear. This question tests not which side you take but how many fears you can see at once.
